% TOP_PROBLEM: {"id": 30007320, "title": "Halting detection with one additional decay power per time derivative", "category_id": 9, "category": {"id": 9, "name": "pde", "display_name": "Partial Differential Equations", "description": "PDEs and their applications in physics and geometry.", "slug": "pde", "order_index": 9, "created_at": "2026-07-31T15:26:25.671Z"}, "status": "open"}
% FIRST_POSTED: 2026-10-06
% ATTEMPT_STATUS: solved
% OPENAI_CLAIM: {"schema_version":1,"source_repo":"https://github.com/openai/math","source_commit":"adc7f1241b42e322a6451854ab7e4b4c146bf78a","release_date":"2026-10-06","match":"full","resolution":"proved","independently_reviewed":false,"second_pass":"agreed","families":[{"family":"376","title":"Universal computation in forced Navier–Stokes flows","subject":"Partial differential equations","manuscripts":[{"title":"Scalar Potentials and Slow Clocks for Forced Fluid Computation","pdf_path":"preprints/Scalar-Potentials-and-Slow-Clocks-for-Forced-Fluid-Computation-September-27-2026/manuscript.pdf","readme_path":"preprints/Scalar-Potentials-and-Slow-Clocks-for-Forced-Fluid-Computation-September-27-2026/README.md","date":"2026-09-27","theorem_ref":"bcs:prefix-theorem (bcs-prefix.tex:4)","statement_sources":[{"path":"preprints/Scalar-Potentials-and-Slow-Clocks-for-Forced-Fluid-Computation-September-27-2026/build/bcs-prefix.tex","line":4,"label":"bcs:prefix-theorem"}]}],"lean":{"doc_path":"lean/docs/376.md","comparators":[],"declarations":[],"covers_main_theorem":false},"reasoning_trace":null}]}
\documentclass[11pt]{article}
\usepackage[margin=25mm]{geometry}
\usepackage{fontspec}
\setmainfont{Latin Modern Roman}
\usepackage{amsmath,amssymb,mathtools}
\usepackage{unicode-math}
\setmathfont{Latin Modern Math}
\usepackage{enumitem}
\usepackage{xurl,hyperref,bookmark}
\hypersetup{colorlinks=true,urlcolor=blue}
\setcounter{secnumdepth}{0}
\setlength{\parindent}{0pt}
\setlength{\parskip}{5pt}
\setlength{\emergencystretch}{3em}
\begin{document}
\section*{Halting detection with one additional decay power per time derivative}

\subsection{OpenAI's claim}
\textit{Scalar potentials and slow-clock computation}: The paper claims compactly supported forced fluid computation with a fixed particle and half-space detector, with time derivatives decaying at the stated inverse-time rates. Other realizations include invariant-plane instructions and velocity-field detection under a companion diffusion statement.

\subsection{Mathematical statement of the claimed result}
\paragraph{Mathematical claim 1.} For the incompressible system $\partial_tu+(u\cdot\nabla)u=-\nabla p+\nu\Delta u+f$, $\nabla\cdot u=0$, $u(0)=0$, choose viscosity $\nu>0$, whose value is computable. The force is smooth and effectively constructed on $\mathbb T^3$, with spatial mean zero. Its corresponding velocity is unique and globally smooth, and both fields satisfy \[\|\partial_t^jD_x^\alpha f(t)\|_\infty+\|\partial_t^jD_x^\alpha u(t)\|_\infty\le C_{j,\alpha}(1+t)^{-1-j}\quad(j\ge0).\] For initial particle position $(1/4,1/2,1/2)$, the event of entering $\{1/2<x_1<1\}$ is equivalent to halting of the supplied machine. A first force choice has zero pressure; its solenoidal projection gives a divergence-free choice with unchanged velocity and altered normalized pressure. Only loading depends on the input for a fixed machine.

\subsection{Outline of the argument}
\textit{Scalar potentials and slow-clock computation}: Three scalar potentials successively lift a coded rectangle, apply its rule, and lower the image. A logarithmic time change remains onto an unbounded computational-time interval while reducing all time derivatives. Area-changing prefix maps are balanced by a normal component, and the alternative planar processor provides the stated stationary, periodic, and diffusion-based detector variants.

\subsection{Scope relative to this problem}
\textit{Scalar potentials and slow-clock computation}: The slow-clock torus realization retains one fixed particle and strip while adding a decay power for each time derivative. The optional divergence-free projection alters the pressure representation while keeping the same velocity and observation event.

With the hypotheses and conventions above, the claimed result addresses the whole mathematical assertion on this page.

\subsection{Formal verification}
Family 376 describes Lean coverage for the selected sheet-program, rapidly decaying particle-detector, velocity-detector, and full-box transport statements described separately in its Comparator files.

\textit{Scalar potentials and slow-clock computation}: No manuscript-specific Comparator is listed in the pinned inventory.

The documented selection does not cover the entire manuscript assertion relevant here.

This site has not run the Lean development. These links describe OpenAI’s selected formal statements and should be read with their stated scope.

\subsection{Status caveat}
OpenAI reports that an internal model produced the manuscripts. Their arguments have not been independently verified as proofs by this site. Any portion outside the documented formal scope may contain mistakes, and the catalogue label reports the announced mathematical claim.

\subsection{Sources}
\begin{itemize}
\item OpenAI, \emph{Scalar Potentials and Slow Clocks for Forced Fluid Computation}, 2026-09-27; citation key \texttt{OAI:Scalar-Potentials-and-Slow-Clocks-for-Forced-Fluid-Computation-September-27-2026}. \href{https://github.com/openai/math/blob/adc7f1241b42e322a6451854ab7e4b4c146bf78a/preprints/Scalar-Potentials-and-Slow-Clocks-for-Forced-Fluid-Computation-September-27-2026/manuscript.pdf}{PDF (GitHub)}; \href{https://github.com/openai/math/raw/adc7f1241b42e322a6451854ab7e4b4c146bf78a/preprints/Scalar-Potentials-and-Slow-Clocks-for-Forced-Fluid-Computation-September-27-2026/manuscript.pdf}{PDF (direct)}; \href{https://github.com/openai/math/blob/adc7f1241b42e322a6451854ab7e4b4c146bf78a/preprints/Scalar-Potentials-and-Slow-Clocks-for-Forced-Fluid-Computation-September-27-2026/README.md}{Manuscript page and citation}; \href{https://github.com/openai/math/blob/main/preprints/Scalar-Potentials-and-Slow-Clocks-for-Forced-Fluid-Computation-September-27-2026/manuscript.pdf}{Latest version}; Statement source: \href{https://github.com/openai/math/blob/adc7f1241b42e322a6451854ab7e4b4c146bf78a/preprints/Scalar-Potentials-and-Slow-Clocks-for-Forced-Fluid-Computation-September-27-2026/build/bcs-prefix.tex\#L4}{bcs:prefix-theorem}.
\item Family 376: \href{https://github.com/openai/math/blob/adc7f1241b42e322a6451854ab7e4b4c146bf78a/lean/docs/376.md}{Lean scope}.
\item \href{https://github.com/openai/math/blob/adc7f1241b42e322a6451854ab7e4b4c146bf78a/CONTENTS.md}{OpenAI catalogue entry}.
\end{itemize}
\end{document}
